\documentclass[10pt]{beamer}

\usetheme{metropolis}
\definecolor{WiLabRed}{RGB}{197,18,48}
\setbeamercolor{frametitle}{fg=white,bg=WiLabRed}
\setbeamercolor{progress bar}{fg=WiLabRed!90}
\setbeamercolor{title separator}{fg=WiLabRed!90}
\setbeamercolor{progress bar in section page}{fg=WiLabRed!90}
\setbeamercolor{background canvas}{bg=white}
\setbeamercolor{alerted text}{fg=WiLabRed!90}

\usepackage{appendixnumberbeamer}

\usepackage{booktabs}
\usepackage[scale=2]{ccicons}

\usepackage{pgfplots}
\usepgfplotslibrary{dateplot}

\usepackage{xspace}
\newcommand{\themename}{\textbf{\textsc{metropolis}}\xspace}

\usepackage{marvosym}
%\usepackage{subfig}
\usepackage{graphicx}\graphicspath{{images/}}
\usepackage{subcaption}
\usepackage[framed]{./support/mcode}
\usepackage{listings}

\title{Laboratory 0: Understanding SDR Hardware}
\subtitle{\textit{Software Defined Radio for Engineers} (Collins~\textit{et al.}), \textsection~5.1-5.2}
\date{}
\author{\textbf{Alexander M. Wyglinski, Ph.D.}}
\institute{ \vspace*{1in}\hfill\includegraphics[height=1.125cm]{wilab_logo-A70916.eps} \qquad \includegraphics[height=1.125cm]{WPI_Inst_Prim_FulClr.eps}}
% \titlegraphic{\hfill\includegraphics[height=1.5cm]{logo.pdf}}

% Foot for all slides
\setbeamertemplate{frame footer}{\tiny \copyright~2018 by Alexander Wyglinski. This work is licensed under the Creative Commons Attribution-ShareAlike 4.0 International License. To view a copy of this license, visit http://creativecommons.org/licenses/by-sa/4.0/.}

\begin{document}

%\captionsetup[subfigure]{labelformat=empty}

%%%%%%%%%%%%%%%%%%%%%%%%%%%%%%%%%%%%%%%%%%%%%%%%%%%%%%%%%%

\maketitle




%%%%%%%%%%%%%%%%%%%%%%%%%%%%%%%%%%%%%%%%%%%%%%%%%%%%%%%%%%

\begin{frame}[fragile]{What Is A Communication System?}


\centering
Information Source\\
$\downarrow$\\
Conversion to EM Signal\\
$\downarrow$\\
Transmission Medium\\
$\downarrow$\\
Conversion from EM Signal\\
$\downarrow$\\
Information Sink

\end{frame}



%%%%%%%%%%%%%%%%%%%%%%%%%%%%%%%%%%%%%%%%%%%%%%%%%%%%%%%%%%

\begin{frame}[fragile]{SDR Basic Ingredients}


\begin{itemize}
 \item \textbf{Analog RF section}: Antennas, RF filters, LNAs, attenuators, mixers
 \item \textbf{Analog baseband section}: Analog filters, ADCs, DACs
 \item \textbf{Digital signal processing section}: FPGAs, DSPs, ARMs, general purpose microprocessors
\end{itemize}

\end{frame}




%%%%%%%%%%%%%%%%%%%%%%%%%%%%%%%%%%%%%%%%%%%%%%%%%%%%%%%%%%

\begin{frame}[fragile]{Transmit \& Receive Specifications}


\begin{itemize}
 \item 300 - 3800~GHz carrier frequency range
 \item 200 - 20000~kHz channel bandwidth
 \item 65.1 - 61440~kSPS
 \item 2.4~Hz LO step size with 5~Hz sample rate step size
 \item 12-bit DACs
\end{itemize}

\end{frame}




%%%%%%%%%%%%%%%%%%%%%%%%%%%%%%%%%%%%%%%%%%%%%%%%%%%%%%%%%%

\begin{frame}[fragile]{Linux Industrial Input/Output (IIO)}


\begin{itemize}
 \item \textit{Linux kernel IIO driver} runs inside the Linux kernel of the ARM on the Pluto SDR
 \item \textit{libiio} userspace library is used for accessing local and remote IIO devices and is located on both Pluto SDR ARM and host machine
 \item \textit{IIO Daemon} (iiod) allows for remote connection to IIO clients such as ARM inside Pluto SDR
\end{itemize}
\vspace*{5mm}
$\rightarrow$ How does this IIO framework look like?

\end{frame}




%%%%%%%%%%%%%%%%%%%%%%%%%%%%%%%%%%%%%%%%%%%%%%%%%%%%%%%%%%

\begin{frame}[fragile]{MATLAB IIO Client}

\begin{itemize}
 \item We can use MATLAB as an IIO client on the host machine
 \item This is achieved using the following system objects:
 \begin{itemize}
  \item \textit{comm.SDRRxPluto}: Pluto SDR receiver system object
  \item \textit{comm.SDRTxPluto}: Pluto SDR transmitter system object
 \end{itemize}
 \item Instantiating system objects:
\end{itemize}
\begin{figure}
\centering
\begin{minipage}[framed]{0.9\textwidth}
\begin{lstlisting}
rx = sdrrx('Pluto');
tx = sdrtx('Pluto');
\end{lstlisting}
\end{minipage}
\captionsetup{labelformat=empty}
\end{figure}

\end{frame}


%%%%%%%%%%%%%%%%%%%%%%%%%%%%%%%%%%%%%%%%%%%%%%%%%%%%%%%%%%

\begin{frame}[fragile]{\texttt{rx} Specifications}

\begin{figure}
\centering
\begin{minipage}[framed]{0.9\textwidth}
\begin{lstlisting}
>> rx = sdrrx('Pluto')

rx =

  comm.SDRRxPluto with properties:

                DeviceName: 'Pluto'
                   RadioID: 'usb:0'
           CenterFrequency: 2.4000e+09
                GainSource: 'AGC Slow Attack'
            ChannelMapping: 1
        BasebandSampleRate: 1000000
            OutputDataType: 'int16'
           SamplesPerFrame: 3660
    ShowAdvancedProperties: false
\end{lstlisting}
\end{minipage}
\captionsetup{labelformat=empty}
\end{figure}

\end{frame}

%%%%%%%%%%%%%%%%%%%%%%%%%%%%%%%%%%%%%%%%%%%%%%%%%%%%%%%%%%

\begin{frame}[fragile]{\texttt{rx} Specifications}

\begin{itemize}
 \item \textit{RadioID}: Related to interface number of connected device, either via USB or Ethernet/IP address
 \item \textit{CenterFrequency}: Operating center frequency of radio in hertz 
 \item \textit{BasebandSampleRate}: Baseband sampling rate of I/Q chain
 \item \textit{GainSource}: Three types are available
 \begin{itemize}
  \item Manual
  \item AGC Slow Attack
  \item AGC Fast Attack
 \end{itemize}
 \item \textit{ChannelMapping}: Set this value to 1 for Pluto SDR
 \item \textit{OutputDataType}: MATLAB can only receive 16-bit complex integers
 \item \textit{SamplesPerFrame}: Defines number of samples in buffer which is passed from MATLAB to iiod
\end{itemize}

 
\end{frame}


%%%%%%%%%%%%%%%%%%%%%%%%%%%%%%%%%%%%%%%%%%%%%%%%%%%%%%%%%%

\begin{frame}[fragile]{Radio I/O Basics: Collecting Data}

\begin{figure}
\centering
\begin{minipage}[framed]{0.9\textwidth}
\begin{lstlisting}
% Perform data collection then offline processing
data = zeros(frameSize, framesToCollect);
for frame = 1:framesToCollect  
    [d,valid,of] = rx();
    if ~valid
        warning('Data invalid')
    elseif of
        warning('Overflow occurred')
    else
        data(:,frame) = d;
    end
end
% Process new live data
sa1 = dsp.SpectrumAnalyzer;
for frame = 1:framesToCollect
    sa1(data(:,frame)); % Algorithm processing
end
\end{lstlisting}
\end{minipage}
\captionsetup{labelformat=empty}
\end{figure}

 
\end{frame}



%%%%%%%%%%%%%%%%%%%%%%%%%%%%%%%%%%%%%%%%%%%%%%%%%%%%%%%%%%

\begin{frame}[fragile]{Radio I/O Basics: Continuous Transmission}

\begin{figure}
\centering
\begin{minipage}[framed]{0.9\textwidth}
\begin{lstlisting}
% Transmit all zeros
tx = sdrtx('Pluto');
fs = 1e6; fc = 1e4; s = 2*pi*fs*fc*(1:2^14).';
wave = complex(cos(s),sin(s));
tx.transmitRepeat(wave);
\end{lstlisting}
\end{minipage}
\captionsetup{labelformat=empty}
\end{figure}
\vspace*{3mm}
Why is this important?\\
$\rightarrow$~When the Pluto SDR is ``on'', it is still transmitting even if the user is not sending anything at that moment

 
\end{frame}

%%%%%%%%%%%%%%%%%%%%%%%%%%%%%%%%%%%%%%%%%%%%%%%%%%%%%%%%%%


\end{document}
